\documentclass{article}
\usepackage[utf8]{inputenc}
\usepackage[T1]{fontenc}
\usepackage[french]{babel}
\usepackage{booktabs}
\usepackage{tabularx}
\usepackage{geometry}
\geometry{margin=2cm}

\title{Analyse du Cas Pratique : Free}
\author{Maël - BTS SIO}
\date{30 Mars 2026}

\begin{document}

\maketitle

\section{Analyse PESTEL du macro-environnement de Free}

Le tableau suivant classe les facteurs d'influence identifiés dans les documents selon les six dimensions de la méthode PESTEL.

\begin{table}[h]
\small
\begin{tabularx}{\textwidth}{l|X}
\toprule
\textbf{Composante} & \textbf{Éléments identifiés} \\
\midrule
\textbf{Politique} & Rôle de l'Arcep (régulateur des télécoms) qui surveille la qualité des réseaux et le respect des engagements des opérateurs. \\
\midrule
\textbf{Économique} & Inflation persistante poussant les consommateurs à faire des arbitrages budgétaires ; Contexte macroéconomique difficile ; Marché du mobile en ralentissement. \\
\midrule
\textbf{Socio-culturel} & Attente des clients pour des prix bas (forfaits à bas coûts) ; Manque d'appétence initial pour la 5G ; Importance croissante du divertissement (streaming). \\
\midrule
\textbf{Technologique} & Évolution vers le Wi-Fi 7 ; Déploiement de la fibre optique ; Augmentation des débits (8 Gb/s) ; Innovation constante dans les box (Freebox Ultra). \\
\midrule
\textbf{Écologique} & (Non mentionné explicitement dans les extraits, mais implicite via les normes de consommation énergétique des équipements). \\
\midrule
\textbf{Légal} & Mise en demeure par l'Arcep pour la remise en état des réseaux en fibre optique mal conçus. \\
\bottomrule
\end{tabularx}
\end{table}

\section{Opportunités et Menaces du macro-environnement}

\begin{itemize}
    \item \textbf{Opportunités :}
    \begin{itemize}
        \item \textbf{Technologiques :} L'arrivée de nouvelles normes comme le Wi-Fi 7 permet de lancer des produits premium (Freebox Ultra) pour se différencier.
        \item \textbf{Socio-culturelles :} La demande pour des offres de divertissement intégrées (Netflix, Disney+, Canal+) permet de créer des offres à forte valeur ajoutée.
    \end{itemize}
    \item \textbf{Menaces :}
    \begin{itemize}
        \item \textbf{Économiques :} L'inflation réduit le pouvoir d'achat, ce qui peut freiner le recrutement de nouveaux abonnés si les prix augmentent.
        \item \textbf{Légales/Réglementaires :} La pression de l'Arcep sur la qualité de service oblige à des investissements imprévus pour corriger les réseaux défaillants.
        \item \textbf{Concurrentielles (Micro-environnement) :} La guerre des prix sur le mobile (baisse d'un tiers du prix moyen en un an) réduit les marges.
    \end{itemize}
\end{itemize}

\section{Impact sur la situation de Free}

L'impact de ces facteurs oblige Free à adopter une stratégie de \textbf{différenciation par l'innovation} tout en maintenant une \textbf{compétitivité prix}. 

\begin{itemize}
    \item \textbf{Innovation de produit :} Le lancement de la Freebox Ultra répond à la saturation du marché en proposant des débits records et un catalogue de services inédit pour justifier le maintien des tarifs.
    \item \textbf{Performance financière :} Malgré la guerre des prix, Free enregistre une hausse de 10,4 \% de son chiffre d'affaires en 2024, prouvant que son modèle d'« entrepreneur innovateur » (au sens de Schumpeter) fonctionne pour gagner des parts de marché.
    \item \textbf{Qualité de service :} La surveillance de l'Arcep impose à Free de ne plus seulement miser sur le prix ou la technologie, mais aussi sur la fiabilité de ses infrastructures (fibre).
\end{itemize}

\end{document}